\chapter{Einleitung}
\label{chapter:kapitel1}

Dieses Kapitel führt in den thematischen und fachlichen Kontext der Arbeit ein und erläutert
die Problemstellung, die Zielsetzung der Untersuchung, die Methodik, sowie den Aufbau der Arbeit.

% Hier beginnt der Text Ihrer Projektarbeit.
% Empfehlung: Einleitung in Problemstellung, Zielsetzung und Aufbau gliedern.


\section{Problemstellung und Ausgangssituation}
\label{section:problemstellung}
Die Instandhaltung industrieller Anlagen verlagert sich zunehmend von papierbasierten Verfahren hin zu digital gestützten, 
mobilen Lösungen, die eine Dokumentation direkt am Einsatzort ermöglichen.\footcite[Vgl.][S. 291]{eisenmann_von_2023} 
Die zugrunde liegenden Prozesse werden häufig im \ac{SAP}-Modul \ac{PM} (heute auch als SAP \ac{EAM} 
bezeichnet) abgebildet. Für den mobilen Zugriff kommen vermehrt Anwendungen auf Basis der Oberflächentechnologie SAP Fiori zum Einsatz.

Im betrachteten Fall wird eine SAP-Fiori-Anwendung zur Bilddokumentation von Instandhaltungsmeldungen genutzt. Anwender können Bilder 
zu Meldungen erfassen, prüfen und mit dem zugehörigen PM-Datensatz verknüpfen. Die Anwendung wurde bereits produktiv eingesetzt, 
ihre Gebrauchstauglichkeit jedoch bislang nicht systematisch untersucht.

Gebrauchstauglichkeit beschreibt das Ausmaß, in dem ein Produkt, System oder eine Dienstleistung von bestimmten Nutzenden in einem 
definierten Nutzungskontext effektiv, effizient und zufriedenstellend verwendet werden kann.\footcite[Vgl.][Abschnitt~3.1]{DIN_ISO_9241-11:2018}
Im laufenden Betrieb wurden Hinweise auf Bedienprobleme und potenzielle Effizienzverluste festgestellt,\footnote{Vgl. Gespräch mit Kaiser 2026, Anhang~\ref{anhang:gespraech}.} deren Ursachen bislang nicht 
strukturiert analysiert wurden. Dadurch fehlt eine fundierte Grundlage zur Bewertung der Gebrauchstauglichkeit und zur gezielten 
Ableitung von Verbesserungsmaßnahmen.

Als anerkannter Bewertungsmaßstab können die Interaktionsprinzipien der ISO 9241-110 (2020) herangezogen werden. Aus der fehlenden 
systematischen Bewertung der Gebrauchstauglichkeit ergibt sich der Bedarf einer strukturierten Untersuchung, aus dem die nachfolgend 
dargestellte Zielsetzung und Forschungsfrage abgeleitet werden.

\section{Zielsetzung und Forschungsfrage}
\label{section:zielsetzung}
Ziel der vorliegenden Projektarbeit ist es, die Gebrauchstauglichkeit der untersuchten SAP-Fiori-Anwendung zur mobilen Instandhaltung 
systematisch und nachvollziehbar zu bewerten und auf dieser Grundlage konkrete Optimierungspotenziale aufzuzeigen. 
Damit soll die in der Problemstellung dargelegte Lücke einer bislang fehlenden, objektivierbaren Bewertungsgrundlage geschlossen werden.

Die \textbf{Forschungsfrage} lautet: Inwieweit erfüllt eine SAP-Fiori-Anwendung zur mobilen Instandhaltung die Interaktionsprinzipien der ISO 9241-110 
und welche Optimierungspotenziale lassen sich aus einer heuristischen Evaluation ableiten?


\section{Methodik der Arbeit}
\label{section:methodik}

Zur Beantwortung der Forschungsfrage wird ein zweistufiges Vorgehen gewählt, das einen theoretisch-konzeptionellen und einen darauf 
aufbauenden evaluativ-anwendungsorientierten Teil enthält. Im ersten Schritt werden die begrifflichen und normativen Grundlagen zu 
Gebrauchstauglichkeit, den Interaktionsprinzipien der ISO 9241-110 sowie zur Methode der heuristischen Evaluation auf Basis 
entsprechender Fachliteratur und der relevanten Normen aufgearbeitet. Sie bilden den Bezugsrahmen für die anschließende Untersuchung.

Als Untersuchungsmethode dient die heuristische Evaluation in Form einer Experteninspektion. Dabei wird die Anwendung anhand etablierter
Bewertungskriterien systematisch auf Usability-Probleme geprüft, ohne dass nutzerbasierte Tests durchgeführt werden. Die Wahl dieses Verfahrens 
ergibt sich aus seiner ressourcenökonomischen und expertenorientierten Ausrichtung, die eine strukturierte Bewertung der Gebrauchstauglichkeit 
auch ohne umfangreiche Nutzerstudien ermöglicht. In Abgrenzung zu nutzerbasierten Verfahren beruht die Bewertung damit auf einer 
kriteriengeleiteten Inspektion und nicht auf der Beobachtung tatsächlichen Nutzerverhaltens.

Die ISO 9241-11 ordnet ein solches Vorgehen methodisch ein, indem sie zwischen einer Prüfung durch Benutzer und einer expertenbasierten Evaluierung
unterscheidet und die Überprüfung auf Abweichungen von festgelegten Kriterien sowie die Erfassung von Anzahl und Schweregrad der Probleme 
als anerkannte Evaluierungsansätze benennt.\footcite[Vgl.][Abschnitt~B.1]{DIN_ISO_9241-11:2018}

Als Bewertungsmaßstab werden die Interaktionsprinzipien der ISO 9241-110 zu einem Kriterienkatalog operationalisiert, der die normativen 
Vorgaben in überprüfbare Bewertungskriterien überführt und als Soll-Maßstab der Untersuchung dient.\footcite[Vgl.][Abschnitt~1]{DIN_ISO_9241-110:2020} Gegen
diesen Katalog wird die Anwendung systematisch geprüft. Die identifizierten Usability-Probleme werden dokumentiert, nach ihrem Schweregrad 
eingeordnet und abschließend zu konkreten Optimierungspotenzialen ausgearbeitet.

Eine methodische Einschränkung besteht darin, dass der Verfasser die im Rahmen dieser Untersuchung 
identifizierten Abweichungen im Anschluss an die Evaluation selbst als Entwickler beheben wird. 
Diese nachgelagerte Rollenübernahme wird transparent gemacht und im weiteren Verlauf methodisch 
reflektiert (Kapitel~\ref{grenzen}).

\section{Aufbau und Struktur der Arbeit}
\label{section:struktur}
Die Arbeit umfasst sechs Kapitel. In Kapitel \ref{chapter:kapitel2} werden die theoretischen Grundlagen der Heuristischen Evaluation, sowie der 
Interaktionsprinzipien vermittelt. Kapitel \ref{chapter:kapitel3} beschreibt den Anwendungskontext, indem es die Instandhaltungsprozesse und ihre 
Abbildung in SAP Fiori als technologische Grundlage sowie den konkreten Untersuchungsgegenstand und dessen Nutzungskontext darstellt. 
Kapitel \ref{chapter:kapitel4} bildet den anwendungsorientierten Kern der Arbeit, in dem die Interaktionsprinzipien zu einem Kriterienkatalog 
operationalisiert, der Erfüllungsgrad der Anwendung ermittelt, die Schwachstellen nach Schweregrad bewertet und daraus 
Optimierungspotenziale abgeleitet werden. Kapitel \ref{chapter:kapitel5} ordnet die Befunde in den Gesamtkontext ein und reflektiert die methodischen 
Grenzen der Untersuchung. Kapitel \ref{chapter:kapitel6} fasst die wesentlichen Ergebnisse zusammen und schließt mit einer kritischen Würdigung sowie einem 
Ausblick.